% ---------------------------------------------------------------------
%  PHỤ LỤC
% ---------------------------------------------------------------------
\chapter{Phụ lục}
\label{ch:phuluc}
\markboth{Phụ lục}{}

% Bảng/hình trong phụ lục đánh số C.1 thay vì 6.1 để rõ đây là bảng phụ lục
\renewcommand{\thetable}{C.\arabic{table}}
\setcounter{table}{0}

\section*{A. Hướng dẫn cài đặt và chạy chương trình}
\phantomsection
\addcontentsline{toc}{section}{A. Hướng dẫn cài đặt và chạy chương trình}

Các bước dưới đây đủ để hội đồng chạy lại toàn bộ hệ thống sàng lọc sức khỏe
trên một máy Windows sạch, chưa từng cài đồ án.

\begin{lstlisting}[language=bash, caption={Cài đặt và khởi chạy hệ thống sàng lọc sức khỏe.}]
# Clone và cài đặt
 git clone https://github.com/SmallFCraft/QuanLyHieuThuoc-Laravel.git pharmacy-laravel
 cd pharmacy-laravel
 composer install
 npm ci && npm run build

# Cấu hình PowerShell + SQLite mặc định
 Copy-Item .env.example .env
 php artisan key:generate
 New-Item -ItemType File -Force database/database.sqlite
 php artisan migrate --seed

# AI tùy chọn: điền AI\_BASE\_URL / AI\_API\_KEY / AI\_MODEL trong .env
# Không có AI, động cơ luật vẫn phân tầng; Briefing chuyển sang chờ.

# Mở hai terminal để chạy app và worker
 php artisan serve --port=8123
 php artisan queue:work --tries=1 --timeout=120

# Mở http://127.0.0.1:8123/health-check và chạy test
 php artisan test tests/Feature/TriageAiBriefingTest.php
# Xem README gốc để biết đầy đủ cấu hình và bộ kiểm thử.
\end{lstlisting}

\begin{luuy}
Nếu không cấu hình khóa API, hệ thống vẫn phân tầng an toàn bằng động cơ luật
tất định; chỉ mô-đun diễn giải AI Briefing chuyển sang trạng thái chờ.
\end{luuy}

\section*{B. Bảng tự kiểm trước khi nộp}
\phantomsection
\addcontentsline{toc}{section}{B. Bảng tự kiểm trước khi nộp}

\begin{tukiem}
  \item Sản phẩm chạy được trên máy khác, không chỉ trên máy của nhóm.
  \item Đã chạy thử demo sàng lọc trên đúng môi trường sẽ dùng khi bảo vệ.
  \item Mã nguồn nằm trên kho \code{git}, lịch sử commit trải đều tiến độ.
  \item Có tệp \code{README}, \code{composer.json} và \code{package.json} để chạy lại.
  \item Mọi hình, bảng đều có chú thích, có nhãn, và được nhắc tới trong thân bài.
  \item Mọi mục tiêu ở phần Mở đầu đều được đối chiếu ở
        Chương~\ref{ch:ketluan}.
  \item Mọi con số trùng khớp với kết quả chạy thật tại Chương~\ref{ch:ketqua}.
  \item Mọi tài liệu trong \code{refs.bib} đều được trích dẫn ít nhất một lần.
  \item Đã khai báo phần có dùng công cụ AI hỗ trợ và giải thích được toàn bộ hệ thống.
\end{tukiem}

\section*{C. Chi tiết kết quả thực nghiệm 26 ca bệnh mẫu}
\phantomsection
\addcontentsline{toc}{section}{C. Chi tiết kết quả thực nghiệm 26 ca bệnh mẫu}

Bảng~\ref{tab:bench_detail} liệt kê toàn bộ 26 ca bệnh kiểm thử, nhóm mong đợi
và nhóm thực tế do hệ thống đưa ra ở lượt chạy thực nghiệm chính thức. Dữ liệu
thô của cả ba lượt chạy (26 ca kiến trúc lai, 26 ca chỉ luật, 6 ca diễn giải)
được lưu dạng JSON tại thư mục \code{benchmarks/} kèm theo báo cáo, kèm câu
lệnh tái lập: \code{php storage/app/bench\_runner.php} (kiến trúc lai),
\code{AI\_TRIAGE\_FACTS=false php storage/app/bench\_runner.php} (chỉ luật),
\code{php storage/app/brief\_runner.php} (diễn giải).

\begin{table}[H]
\centering
\caption{Chi tiết kết quả phân tầng 26 ca thực nghiệm}
\label{tab:bench_detail}
\small
\setlength{\tabcolsep}{4.5pt}
\begin{tabular}{|l|p{4.2cm}|l|l|l|}
\toprule
\textbf{Mã ca} & \textbf{Bệnh cảnh} & \textbf{Mong đợi} & \textbf{Thực tế} & \textbf{Đánh giá} \\
\midrule
EM\_01 & Tim mạch & Cấp cứu & Cấp cứu & Đúng \\
EM\_02 & Đột quỵ & Cấp cứu & Dược sĩ & Thiếu hụt \\
EM\_03 & Nhi khoa & Cấp cứu & Cấp cứu & Đúng \\
EM\_04 & Xuất huyết tiêu hoá & Cấp cứu & Cấp cứu & Đúng \\
EM\_05 & Phản vệ & Cấp cứu & Cấp cứu & Đúng \\
EM\_06 & Rối loạn ý thức & Cấp cứu & Cấp cứu & Đúng \\
\midrule
RD\_01 & Nhiễm trùng dai dẳng & Đi khám & Đi khám & Đúng \\
RD\_02 & Nhi khoa dưới 1 tuổi & Đi khám & Đi khám & Đúng \\
RD\_03 & Mất nước & Đi khám & Tự chăm sóc & Thiếu hụt \\
RD\_04 & Thần kinh kéo dài & Đi khám & Cấp cứu & Bảo thủ \\
RD\_05 & Sản phụ khoa & Đi khám & Đi khám & Đúng \\
\midrule
PR\_01 & Thuốc không đáp ứng & Dược sĩ & Dược sĩ & Đúng \\
PR\_02 & Thai kỳ & Dược sĩ & Đi khám & Bảo thủ \\
PR\_03 & Cho con bú & Dược sĩ & Đi khám & Bảo thủ \\
PR\_04 & Bệnh nền & Dược sĩ & Tự chăm sóc & Thiếu hụt \\
PR\_05 & Không đặc hiệu & Dược sĩ & Dược sĩ & Đúng \\
\midrule
SC\_01 & Đau đầu căng thẳng & Tự chăm sóc & Tự chăm sóc & Đúng \\
SC\_02 & Viêm họng nhẹ & Tự chăm sóc & Tự chăm sóc & Đúng \\
SC\_03 & Cảm lạnh & Tự chăm sóc & Tự chăm sóc & Đúng \\
SC\_04 & Khó tiêu & Tự chăm sóc & Tự chăm sóc & Đúng \\
SC\_05 & Đau cơ & Tự chăm sóc & Dược sĩ & Bảo thủ \\
SC\_06 & Ợ chua & Tự chăm sóc & Tự chăm sóc & Đúng \\
SC\_07 & Khó ngủ & Tự chăm sóc & Tự chăm sóc & Đúng \\
SC\_08 & Đau răng & Tự chăm sóc & Đi khám & Bảo thủ \\
SC\_09 & Ngứa da & Tự chăm sóc & Tự chăm sóc & Đúng \\
SC\_10 & Mỏi khớp & Tự chăm sóc & Tự chăm sóc & Đúng \\
\bottomrule
\end{tabular}
\end{table}

\section*{D. Nhận xét của giảng viên hướng dẫn}
\phantomsection
\addcontentsline{toc}{section}{D. Nhận xét của giảng viên hướng dẫn}

\khungnhanxet[7.5cm]

\vspace{0.8em}
\noindent\begin{minipage}[t]{0.5\textwidth}
\color{daunavy}
{\color{daumaroon}\bfseries Điểm số}

\vspace{0.7em}
\dongchamnhan{Bằng số:}

\vspace{0.9em}
\dongchamnhan{Bằng chữ:}
\end{minipage}\hfill
\okyten{\dauGVHD}{Giảng viên hướng dẫn}
